% 第 45 章　测量、概率与所谓“坍缩”（完整样章）
\chapter{测量、概率与所谓“坍缩”}

第~43~章给了我们一台预测机器：写出量子态，算出振幅，取模的平方，得到每种测量结果的概率。第~44~章告诉我们这台机器在两次测量之间如何运转：薛定谔方程决定状态随时间连续地变化。但有一个问题一直被我们故意跳过——测量那一刻到底发生了什么？教科书里那个吓人的词“坍缩”，说的是同一件事：测量之前是好多可能性，测量之后只剩一个结果。这一步是物理过程，还是记账规则？需要一个人“看”才算数吗？本章的任务是把这个黑匣子拆开，并且把里面的东西分成三堆：哪些是实验反复验证过的事实，哪些是数学模型的计算规则，哪些至今仍是物理学家争论不休的解释故事。学完这一章，你不一定能回答“坍缩到底是什么”，但你应该能一眼看穿关于它的九成流行误读。

\step{反常识开场}

拿出你的手机，把屏幕点亮，再戴上一副偏振太阳镜（很多司机镜、钓鱼镜都是偏振镜）。隔着镜片看屏幕，慢慢转动手机或者镜片：转到某个角度，屏幕会完全变黑。这不奇怪——液晶屏发出的光是偏振光，你可以粗略地想象成光的“抖动方向”被整理到了同一个方向上，而偏振镜片像一排栅栏，只放行顺着缝隙方向的抖动。镜片转到“栅栏与抖动垂直”的位置，光自然全被挡住。

真正奇怪的事情在后面。保持镜片和手机的角度不动，让屏幕保持全黑。现在再找一片偏振片——另一副太阳镜上拆下来的镜片也行——把它斜着插到屏幕和第一片镜片之间，角度大约取四十五度。按照常识，过滤器只会越加越暗：两片十字交叉的栅栏已经挡得一滴不漏，中间再塞一片斜的栅栏，怎么可能更亮？

可它就是更亮了。屏幕上一块明显的亮光重新出现，转动中间那片，亮度还会跟着变化。更奇怪的是：把这片偏振片拿到两片之外——放到最前面或者最后面——什么也不会发生，屏幕依然是黑的。只有夹在中间，它才起作用。

一片“过滤器”，放在某些位置挡光，放在另一些位置反而“放光”。这不是滤镜质量问题，而是量子力学最深处的秘密正躺在你的手机屏幕和两副太阳镜之间。本章结束时，你会能完整解释它，而且会发现：它跟我们如何认识世界、知识如何更新、现实在多大程度上“先于测量而存在”，全是同一个问题。

\step{先猜一猜}

在继续读之前，请把你的预测写下来，写得越具体越好：

两片偏振片已经十字交叉、完全不透光。在中间插入第三片（斜四十五度），你猜屏幕会（a）依然全黑，（b）变亮一点点，（c）明显变亮？如果你猜了（b）或（c），请回答一个更尖锐的问题：第三片挡掉了一部分光，这部分光是“本来就能到达眼睛的”还是“本来就被前两片判了死刑的”？如果是后者，一个过滤器凭什么能“复活”已经不存在的光？

再猜第二个问题。假设我们把光调得极弱极弱，弱到平均每次只有一颗光子飞向这三片偏振片。一颗光子不可再分，它要么整个通过，要么整个被吞掉。你猜：十字交叉的两片之间，一颗光子的命运是“必死”还是“有一半机会活下来”？插进中间那片之后呢？

最后猜一个与光无关的问题。你明天要体检，护士要量你的血压。你猜：把袖带绑上胳膊这个动作，会不会改变你的血压数值？改变得多不多？这个问题听起来和偏振片风马牛不相及，但请记住它——本章后半部分我们会发现，“测量是否打扰被测对象”正是隔开日常世界和量子世界的那道门槛。

写好了吗？把你的答案夹在书页里。读完本章再翻回来对一对，看看自己哪些直觉是对的，哪些直觉需要拆掉重装。

\step{动手实验}

\begin{experiment}[三片偏振片：会“放光”的过滤器]
\textbf{材料}：一部液晶屏手机或电脑显示器（出射光自带偏振）；一副偏振太阳镜；再加一片偏振片——另一副太阳镜拆一片镜片、偏光 3D 眼镜、或淘汰计算器液晶屏前的偏振膜都可以。

\textbf{第一步}：屏幕点亮一片纯白画面（打开备忘录即可）。透过一片偏振片看屏幕，慢慢旋转它，找到屏幕最黑的角度。这说明屏幕光的偏振方向与镜片此时“放行”的方向垂直。

\textbf{第二步}：保持这片不动，把第二片偏振片叠上去并旋转，再次找到全黑。现在两片“栅栏”十字交叉，光被彻底关死。

\textbf{第三步}：保持两片都不动，把第三片斜着插到它们中间，慢慢旋转。观察：亮度从全黑升起，转到约四十五度时最亮。记下最亮时第三片的角度。

\textbf{第四步}：把第三片抽出来，改放到两片的最前面或最后面，再怎么旋转，屏幕始终是黑的。

\textbf{记录与思考}：同一片偏振片，放在中间能“救活”光，放在外面毫无作用。请写下你的解释，尤其回答：中间那片到底对光做了什么？它如果只是“挡掉一部分”，为什么放错位置就完全无效？
\end{experiment}

如果手边实在凑不齐三片，两片也能完成前四步中的大部分观察；甚至只用一副太阳镜对着屏幕旋转，你就能确认“屏幕光是偏振光”这件事本身——很多人戴了多年偏振镜，从未发现自己每天都在和一个量子现象打交道。

\step{旧模型遇到麻烦}

先用我们已有的最好的经典模型来硬算这笔账。把光当作绳波：抖动方向沿着绳子，栅栏只放行顺着缝隙的分量。第一片竖直栅栏把屏幕光整理成竖直抖动；第二片水平栅栏对竖直抖动“全部咬住”，一滴不放。到现在为止一切正常。麻烦出在第三片：它是斜四十五度的栅栏，按照绳波图像，它只放行斜向分量——可此时光里已经没有任何分量了，竖直抖动被第二片吃得干干净净。一个已经归零的东西，无论再过滤多少次都只能是零。绳波模型斩钉截铁地预言：插入第三片，结果依然全黑。实验说：亮了。模型第一次不及格。

那就投降，改说“第三片不是过滤器，它扭转了抖动方向”？这个说法在经典波动里确实能蒙混过去（斜栅栏放行的斜向分量里含有水平分量）。但把光调到极弱，一颗一颗地发射光子，真正的麻烦才登场。实验事实是：十字交叉的两片之间，每颗光子百分之百被吞掉，从没有一颗到达探测器；插入中间那片之后，却有相当比例的光子完好地出现在屏幕上。请注意这有多离谱：对单颗光子而言，“被吞掉”和“到达”之间没有中间态。第二片栅栏本来判决每颗光子“必死”，第三片的到来却把一部分死刑犯变成了幸存者。过滤器不可能复活已经消失的东西——除非“被第一片整理成竖直偏振”这件事本身，并没有给光子刻上一张不可更改的终身档案，而中间那片做的不只是筛选，它还重新改写了光子的状态。

再看血压那道猜题。袖带充气压迫血管，确实会轻微改变血流，但改变量远小于血压计的分辨率；体温计从舌头吸走一点热量，对体温的影响小到完全可以忽略。整个经典物理学建立在一个近乎不言自明的信念上：测量只是“读出”系统本来就有的属性，仪器造成的打扰想要多小就能做多小，原则上可以趋于零。这个信念在日常生活中战无不胜，于是我们不自觉地把它升级成了关于世界的公理：任何对象在任何时刻都同时拥有确定的位置、速度、偏振方向，测量只是把答案抄出来。量子力学恰恰在这条公理上翻了船。如果光子每时每刻都“本来就有”一个确定的偏振方向，那么第二片读出“水平方向：没有”之后，第三片无论读什么，都不该凭空造出水平分量来。唯一的出路是：承认“抄答案”这幅图景在微观世界失效了，测量本身就是一次真实的物理碰撞，它不只是揭示状态，它还会更换状态。

还有第二处翻船现场，比偏振片更刺眼。第~43~章的双缝实验里，一颗一颗发射的电子在屏幕上织出干涉条纹。现在做个改动：在两条缝旁各放一个探测器，记录每颗电子走了哪条缝——条纹立刻消失，屏幕上只剩两团朴素的小山包。有人会说：是探测器撞了电子，把它撞乱了。但实验可以做得很“温柔”：用极其微弱的光去“瞥”电子的行踪，瞥得越轻，条纹只是变淡而不是消失，瞥的程度和条纹模糊的程度严格同步。决定条纹生死的不是那一撞有多重，而是“哪条路”这份信息是否在这个世界上存在过。哪怕没有任何人去读探测器的记录，只要记录原则上可以被读出来，干涉就死了。信息本身就是物理作用——这是旧账本无论如何也记不下的一笔。

顺带把“一颗一颗来”算成实数，免得它听起来像修辞。一支一毫瓦的红色激光笔，每秒发出约 $3\times10^{15}$ 颗光子（每颗红光光子能量约 $3\times10^{-19}$~J）。想让平均每秒只剩一颗，得把光衰减十五个数量级。这听起来离谱，但实验室里用衰减片和单光子源天天在做，逐粒子的双缝、逐粒子的偏振实验早已是常规操作。所以上面说的每个“单颗光子的命运”，都不是思想游戏的虚构，而是有收据的实验事实。

\step{发明新概念}

既然旧的记账方式破产了，我们像历史上的物理学家一样，把账目重新设计一遍。新账簿只有四栏。

第一栏：\textbf{可观测量}。每次测量都是向系统提一个具体的问题，而且问题由你选定。“你的偏振沿竖直方向吗”是一个问题；“你的偏振沿四十五度方向吗”是另一个问题。它们问的是同一个光子，却不是同一笔账——这一点马上会变得生死攸关。日常生活里我们其实也隐约知道提问方式的重要：问卷里“你支持这项改革吗”和“你能容忍这项改革的代价吗”问的是同一批人，统计结果却可能南辕北辙。但那只是心理与措辞的把戏，换种问法之前，受访者心里的答案并不因此改变。量子世界的“提问”要深刻得多：换一个问题，等于换一组物理上完全不同的相互作用。

第二栏：\textbf{可能结果}。对同一个问题，量子理论只提供一张有限或无穷的备选答案清单。偏振提问的清单很短：通过，或者挡住。清单之外没有“半个通过”这种选项——单颗光子的结局总是整份的，概率只在你重复提问许多次之后才浮出水面。

第三栏：\textbf{概率}。量子态（第~43~章的波函数、振幅）配上玻恩规则，给出清单上每个答案出现的概率：概率等于对应振幅的模的平方。理论对单次实验不作承诺，只承诺大量重复之后的比例。这是实验事实：一百年前人们还不愿意接受它，今天它已被数不尽的实验压成了铁案——对完全相同的制备重复测量，结果出现的频率精确地服从玻恩规则。

第四栏，也是本章真正的主角：\textbf{状态更新}。一旦测量给出了某个答案，系统此后的行为就不再由旧态决定，而由“与该答案相容的新态”决定。光子通过了四十五度偏振片之后，它接下来的所有概率都必须从“四十五度偏振态”重新算起——旧态里“它曾经是竖直偏振”的信息被这次相互作用彻底冲销了。这套“测得什么结果，就把态换成对应的新态”的规则，就是所谓坍缩的全部操作内容。请注意：账簿里没有任何一栏提到“有人看见”。探测器的一次电流脉冲、照相底片上一个变黑的颗粒、空气分子被光子撞了一下后留下的轨迹——只要相互作用在世界里留下了不可抹去的记录，状态更新就发生了。意识在这本账里连脚注都不是。

这套规则听起来很怪，但它有一个你从小就熟悉的近亲：贝叶斯更新。抽屉里有一双袜子，一黑一白，你闭着眼摸出一只，是黑的。那么“剩下那只是白的”这条概率立刻从二分之一跳到了一。你的知识被新证据刷新了——公式结构上，这和量子测量后的状态更新一模一样：拿到结果，删去与结果不相容的可能性，剩下的重新归一。但这个比喻必须立刻标明边界。它像的地方是：都是“证据来了，重算概率”。它不像的地方有两个。其一，袜子的颜色在你摸它之前就写在袜子上，贝叶斯更新只改变了你的知识，没改变袜子；而光子通过四十五度偏振片后，是系统本身被推进了新状态——之后的一切实验都按新态行事，旧态不存在任何“档案复印件”可供复查。其二，对袜子你可以先问颜色再问质地，顺序无所谓；对量子系统，先问竖直再问水平，和先问水平再问竖直，是两笔不同的账，结果分布都不一样。你的开场实验已经亲手验证了这一点：第三片只有“插在中间”才管用，说明提问的顺序本身是物理的一部分。至于“有没有一张预先写好所有问题答案的总表”——这个问题太重要，也太反直觉，我们要留到第~48~章的贝尔不等式再正式开庭。这里只先备案：量子力学的预言和“存在这样一张总表”的假设，在数学上不能同时成立。

\step{一句话模型}

\begin{oneline}
量子测量不是偷看一张早已写好的答案表，而是一次真实的物理相互作用：你选定一个问题，大自然按玻恩规则随机给出答案，并把系统推进与这个答案相容的新状态——此后的全部预言，都必须从这张新答卷重新算起。
\end{oneline}

这句话像什么？像洗牌：每测一次，牌堆被真正洗过，之前的顺序作废。它不像什么？不像翻牌：牌面不是在翻开之前就固定朝下等着你的。也请注意它“不是”什么：它不是“仪器太粗糙所以测不准”的道歉信，不是“人的意识创造了现实”的许可证，也不是“一切都不确定、科学破产了”的投降书——恰恰相反，它给出的是自然界最精确的概率预言。

\step{数学翻译器}

现在把开场实验从头翻译成数字。设屏幕光的偏振方向为 $0^\circ$。第~43~章说过，量子态可以按任意一组互相垂直的方向分解，振幅就是分解系数。一个 $0^\circ$ 偏振的光子，相对“$45^\circ$ 方向”这组新问题，分解为
\[
|\,0^\circ\,\rangle=\cos 45^\circ\,|\,45^\circ\,\rangle+\sin 45^\circ\,|\,135^\circ\,\rangle .
\]
这条式子回答四个问题。它描述什么：同一个偏振态，用另一组问题来记账时的写法。为什么这样写：偏振方向是按几何角合成的，分解系数就是投影，即余弦和正弦——和第~4~章向量沿斜面分解用的是同一套几何。何时成立：只要是理想偏振片、理想单色偏振光。何时失效：换成会吸收、会散射的真实镜片时，它只是极好的近似。

玻恩规则立刻给出概率：光子通过 $45^\circ$ 片的概率是 $\cos^2 45^\circ=\tfrac12$。通过之后，状态更新为 $|\,45^\circ\,\rangle$——旧账结清。现在第二问来了：这片 $45^\circ$ 偏振的光子遇到 $90^\circ$ 偏振片，同样分解，通过概率又是 $\cos^2 45^\circ=\tfrac12$。两步合起来，一颗屏幕光子连闯三片的概率是
\[
P=\frac12\times\frac12=\frac14 .
\]
做一遍极限检查：把中间那片抽掉，相当于第二问变成“$0^\circ$ 的光子过 $90^\circ$ 片”，概率 $\cos^2 90^\circ=0$，全黑——对。把中间那片转到与第一片平行，$\cos^2 0^\circ=1$ 乘 $\cos^2 90^\circ=0$，还是全黑——对，因为你的实验里只有斜着才亮。公式在两个极端角度上都通过了盘问，才配得上中间情形的成功。

同一条规则还能直接解释你实验第三步里“慢慢旋转、亮度慢慢变化”的曲线。设入射偏振方向与偏振片放行方向夹角为 $\theta$，单颗光子的通过概率是 $\cos^2\theta$，于是大量光子的通过率——也就是亮度——正比于 $\cos^2\theta$。这正是波动光学里以实验定律身份流传了两百年的马吕斯定律；在量子账簿里，它不再是单独的经验规律，而是玻恩规则对偏振的直接推论。检查一下三个特殊点：$\theta=0^\circ$ 全通过，$\theta=90^\circ$ 全挡住，$\theta=45^\circ$ 恰过一半——和你的镜片完全一致。一条两百年前的宏观定律被还原成单颗光子的概率规则，这种“旧公式在新模型里找到娘家”的感觉，是物理学里少有的幸福时刻。

\begin{figure}[htbp]
\centering
\begin{tikzpicture}[>=Stealth,scale=1.0]
  \draw[->,very thick] (-0.6,0) -- (8.2,0) node[right]{光};
  \foreach \x/\a in {1.2/0, 3.8/45, 6.4/90}{
    \draw[thick,fill=gray!25] (\x,-0.85) rectangle ++(0.18,1.7);
    \node[below] at (\x+0.09,-0.9) {$\a^\circ$ 片};
  }
  \node[above] at (0.3,0.15) {屏幕光};
  \node[above] at (2.5,0.15) {$\tfrac12$ 通过};
  \node[above] at (5.1,0.15) {$\tfrac12$ 通过};
  \node[above] at (7.3,0.15) {$\tfrac14$ 到达};
\end{tikzpicture}
\caption{三片偏振片的量子账本：每一片是一次提问，通过后状态被更新，概率逐级相乘。抽掉中间那片，总概率立刻从 $\tfrac14$ 跌回 $0$。}
\end{figure}

\begin{mathbridge}
\textbf{条件概率与重新归一。}状态更新在数学上朴素得近乎乏味：测得某个结果后，把态里与该结果不相容的分量划掉，剩下的部分重新归一化。设某态按两个可能结果写成 $|\psi\rangle=a\,|\text{通过}\rangle+b\,|\text{挡住}\rangle$，且 $|a|^2+|b|^2=1$。测得“通过”后，新态就是 $|\text{通过}\rangle$ 本身——划掉 $b$ 项，$a$ 项因为自己已经是单位长度，连除法都省了。概率的乘法规则也随之现身：先过 $45^\circ$ 片、再过 $90^\circ$ 片的总概率等于两步概率相乘，这正是条件概率公式 $P(A\,\text{且}\,B)=P(A)\,P(B\mid A)$。量子力学在这里没有发明新的概率论，它发明的是概率\textbf{从哪里来}：不是从你对花色的无知里来，而是从振幅的模平方里来。

\textbf{不确定关系是宽度的关系。}第~44~章说过，一个量子态按位置记账和按动量记账是同一份信息的两种写法，波包压得越窄，动量展开得越宽。把这件事写成不等式，就是海森堡不确定关系
\[
\Delta x\,\Delta p\geq \frac{\hbar}{2},\qquad \hbar\approx 1.05\times10^{-34}\ \text{J·s}.
\]
请务必读准它的主语：$\Delta x$ 和 $\Delta p$ 不是某一次测量的误差，而是\textbf{同一批完全相同地制备出来的系统}，分别测位置和测动量时，结果统计分布的宽度。仪器做得再完美，这两个宽度的乘积也压不到 $\hbar/2$ 以下——窄波包这种“状态本身”就不携带比它更确定的动量信息。这不是工程的限制，是“态”这个概念的内置属性。
\end{mathbridge}

来算一笔真账，看看这条规则把世界分成了两半。取一粒直径约十微米的尘埃，质量约 $5\times10^{-13}$~kg，假设我们把它的位置确定到 $\Delta x=10^{-6}$~m（头发丝直径的量级），那么
\[
\Delta v\geq\frac{\hbar}{2m\,\Delta x}\approx\frac{1.05\times10^{-34}}{2\times5\times10^{-13}\times10^{-6}}\approx10^{-16}\ \text{m/s}.
\]
这个速度不确定度小到什么程度？按它走上一年，累积的偏差不到一颗原子核的直径。尘埃可以同时“有”位置和速度，看不出任何量子痕迹。再看氢原子里的电子：质量 $9.1\times10^{-31}$~kg，位置被原子大小限制在 $\Delta x\approx10^{-10}$~m，于是
\[
\Delta v\geq\frac{\hbar}{2m\,\Delta x}\approx6\times10^{5}\ \text{m/s},
\]
这已经和原子中电子特征速度（约 $10^6$~m/s）同量级——对电子来说，“它此刻在哪里、正以多大速度往哪走”这个经典问题，根本没有足够好的答案可供测量。同一条不等式，除以悬殊的质量，解释了为什么你我一生的直觉都站在经典那边，也解释了为什么那套直觉不能外推到原子里面去。

\begin{challenge}
\textbf{顺序为什么有数学后果。}把“问 $0^\circ$”“问 $45^\circ$”“问 $90^\circ$”记为三种测量 $A$、$B$、$C$。你的实验展示了 $A$、$C$ 连做（十字交叉）必得“挡住”，而 $A$、$B$、$C$ 顺次做却有 $\tfrac14$ 通过。也就是说，$B$ 插在中间改变了结局：对这组测量，先做谁后做谁不是同一件事。第~46~章会把每种测量翻译成一个算符，那时你会发现“顺序要紧”这件事浓缩成一句代数：两个算符的乘积不满足交换律，$AB\neq BA$，而它们的差（对易子）非零的程度，恰好管着不确定关系的下限。不确定性与测量顺序，原来是一枚硬币的两面。

\textbf{退相干：环境是最勤快的观察者。}为什么桌上这粒尘埃从不表演“同时在两处”的叠加？不是因为它大这个抽象属性，而是因为它从不孤单：在普通空气里，它每秒被约 $10^{18}$ 到 $10^{20}$ 个分子撞击，还不断反射、辐射着光子。每一次碰撞都相当于环境对它做了一次“位置测量”，把相位信息撒进无数分子的运动里，再也收不回来。定量估算（约希与策等人的经典结果）表明：一粒十微米的尘埃若被制备成间距一厘米的叠加，其干涉能力在空气中大约 $10^{-30}$ 秒内就消失殆尽——比任何仪器能分辨的时间都短几十个数量级。这就是宏观世界“看起来经典”的动力学原因：不是经典定律接管了，而是环境以不可阻挡的效率持续地“测量”着一切，叠加的相位被稀释成了事实上的经典概率混合。
\end{challenge}

\step{模型的边界}

现在做本章最重要的分层练习：把“坍缩”拆成实验事实、数学模型和物理解释三堆。

\textbf{实验事实}：对相同制备的重复测量服从玻恩规则；某些测量序列的结果依赖顺序；单次结果不可预言、统计分布极其精确；测量需要物理相互作用，无人旁观照样发生（云雾室里的径迹、底片上的颗粒从不需要观众）。这一堆，一百年内没有被任何实验撼动过。

\textbf{数学模型}：用态矢量记账，测量时按玻恩规则取概率、按“划掉不相容分量再归一”更新状态；更一般、更不理想的测量（半透明的探测器、间接读数）有对应的推广语言（POVM），本章只用了其中最理想的投影测量。这套规则何时成立：凡量子力学适用之处，它从未失手。何时可能失效或不够用：它没有回答“为什么这一次得到的是这个结果而不是那个”，也没有说清“相互作用”和“测量”的分界线画在哪里——这就是著名的测量问题，一个真实的、尚未关闭的理论缺口。

\textbf{物理解释}：对同一个数学模型，物理学家讲着好几种不同的故事，实验至今无法区分它们的日常预言，所以任何人把其中一种当成“实验已证实的事实”推销给你，你都有权保持警惕。这里只画一张中立地图。哥本哈根诠释：把坍缩当作获得不可逆记录时的规则更新，量子与经典之间需要一条实用分界线，它好用但拒绝告诉你线画在哪。多世界诠释：干脆删掉坍缩这条规则，只留薛定谔方程，测量使宇宙（包括你）按不同结果分支；“概率在分支里意味着什么”是它的软肋。客观坍缩理论：给方程加上微小的随机自发坍缩项，让大物体自动定域——它不再只是诠释，因为它给出了与标准量子力学略有差别的预言，正被实验一点点挤压。玻姆理论：粒子始终有确定位置，由波函数引导，代价是方程里明目张胆的非局域性；它的预言与标准量子力学相同。QBism：量子态不是世界的属性，而是观测者对自身未来经验的信念度，坍缩就是贝叶斯更新——它把本章的袜子比喻推到极致，同时坦然承认顺序与“无预先答案表”带来的全部后果。五种故事，一台预测机器。区分“共同预言”与“诠释差异”，是阅读一切量子科普时最值钱的防身术。

还有三条边界必须划清。第一，退相干不等于解决问题。退相干漂亮地解释了干涉为何消失、指针为何有确定指向——它把“古怪的叠加”稀释成“看起来古典的各居其一”；但它给出的是一份所有选项仍并列其中的清单，至于“为什么最终落到这其中一个”，它是否算单独给出了答案，至今仍有争议。第二，不确定关系与仪器粗糙无关：再完美的仪器面对的也是同一个下限，因为宽度属于状态本身。第三，坍缩不是超光速开关：对一个粒子的测量更新的是这个粒子的态，至于两个纠缠粒子的情形为什么依然不能用来传信，要等第~48~章细算。

\begin{figure}[htbp]
\centering
\begin{tikzpicture}[>=Stealth,scale=1.0]
  % 左：测量前
  \draw[->] (-0.2,0) -- (2.6,0);
  \draw[->] (0,-0.2) -- (0,1.8);
  \draw[fill=blue!40] (0.5,0) rectangle ++(0.6,1.2);
  \draw[fill=blue!40] (1.5,0) rectangle ++(0.6,1.2);
  \node[below] at (0.8,0) {通过};
  \node[below] at (1.8,0) {挡住};
  \node[above] at (1.3,1.9) {测量前：各 $\tfrac12$};
  % 箭头
  \draw[->,very thick] (3.0,0.8) -- (4.2,0.8) node[midway,above]{测得“通过”};
  % 右：测量后
  \draw[->] (4.6,0) -- (7.4,0);
  \draw[->] (4.8,-0.2) -- (4.8,1.8);
  \draw[fill=orange!60] (5.3,0) rectangle ++(0.6,1.2);
  \draw[fill=gray!20] (6.3,0) rectangle ++(0.6,0.03);
  \node[below] at (5.6,0) {通过};
  \node[below] at (6.6,0) {挡住};
  \node[above] at (6.1,1.9) {测量后：$1$ 与 $0$};
\end{tikzpicture}
\caption{状态更新：测得结果后，不相容的可能性被划掉，剩下的重新归一。形式上与贝叶斯更新相同，物理上却是系统本身换了状态。}
\end{figure}

\step{回头解释开场现象}

回到那片会“放光”的过滤器，现在每一步都有了正式的账目。屏幕光先被整理成 $0^\circ$ 偏振。第一片太阳镜在“$90^\circ$ 问题”上对它提问：$0^\circ$ 与 $90^\circ$ 恰好是一组互斥答案，通过概率 $\cos^2 90^\circ=0$——每颗光子都被吞下，屏幕全黑。注意此刻没有任何东西“藏着”一份水平偏振：账上写得清清楚楚，这个态对这个问题只有“挡住”一个答案。

插入 $45^\circ$ 片，等于在两问之间插入了一个新的提问。对“$45^\circ$ 吗”这个问题，$0^\circ$ 态的答案不再确定：振幅是 $\cos 45^\circ$，一半光子回答“通过”。关键一步是状态更新：回答“通过”的光子被这次相互作用推进了 $45^\circ$ 偏振态，旧档案作废。于是最后一片面对的不再是“必挡”的 $0^\circ$ 光子，而是对 $90^\circ$ 问题有一半通过率的新光子。总通过率 $\tfrac14$，屏幕重新亮起。反直觉的核心一句话说完：我们多加的不是一层滤网，而是一次改写状态的测量——过滤器没有复活死去的光，它改变了“死”的定义所依据的账簿。

第三步实验里“放错位置就无效”也立刻清楚了：把第三片放到两片之外，提问顺序仍是“$0^\circ$ 后直接问 $90^\circ$”，答案照旧是全挡。顺序即物理——这正是挑战盒子里“测量不可交换”的家用版本，也是第~47~章斯特恩—盖拉赫实验将以更锋利的方式再次上演的故事。

这幅图景还有一个分量极重的工程应用：量子保密通信。设想双方用偏振光子逐个传送比特，发送方随机选用两组互不相容的提问方式编码。任何中途截获光子的窃听者都面临死局：她必须选一个方向提问，选对的概率只有一半；一旦选错，她的测量就把光子推进了错误的新态，再转发出去也抹不掉痕迹。合法双方事后公开核对一小部分比特，窃听留下的印记会以约四分之一的错误率赫然显现。这里没有用到任何高深设备，用的正是本章的规则：测量必然改写状态，而改写必然留下统计学脚印。窃听者不是输在仪器不够好，而是输在自然定律不允许“只读不写”的量子测量。

\step{脑洞实验与挑战题}

\begin{thought}[把猫关进盒子，把月亮推过双缝]
薛定谔当年设计那只著名的猫，本意不是宣称猫会半死半活，而是给测量问题拍一张 X 光片：如果微观叠加可以沿“原子—探测器—毒气—猫”这条链条无限放大，那么我们岂不得承认宏观世界也有叠加？本章给了你两层回应。第一层是动力学的：猫每秒与天文数字的分子、光子相互作用，退相干把“活且死”的相位信息在无法想象的短时间内冲散，猫的叠加在实践上等同于“非活即死”的经典清单——没有人能用干涉仪检验一只猫的两种结局之间的相位。第二层是概念上的：退相干解释了“看起来古典”，却没有规定“为什么是这一个结局”，测量问题在猫的尺度上依然开着一道缝。再极端一点：想让月亮表演双缝干涉吗？原则上只需要保证没有任何东西——没有一个光子、一粒尘埃、一次引力扰动——“记下”月亮走了哪条缝。算一算退相干时间，你会发现宇宙本身就是一台永不停歇的测量仪器。“月亮在没人看时是否仍在那里”之所以是个玩笑，是因为月亮从来不缺“看”它的东西。
\end{thought}

\begin{puzzles}
\item 把 $90^\circ$ 的转角切成小步：用四片偏振片依次转 $0^\circ,30^\circ,60^\circ,90^\circ$，总通过率是多少？换成 $n$ 片、每片只转 $90^\circ/n$，极限是多少？这说明了“过滤器”和“测量”的什么差别？\emph{提示}：每片通过概率为 $\cos^2\theta$，总概率是连乘；$n$ 很大时 $\cos^2(90^\circ/n)\approx1-(90^\circ/n)^2$ 换算成弧度再连乘，想想指数趋近于几。结论出乎意料：足够多次“温柔的提问”几乎不挡光。
\item 有人宣布发明了“零打扰显微镜”：看一颗电子而不碰它，从而同时精确测定位置和动量，绕开不确定关系。不用任何公式，指出这个方案的致命伤。\emph{提示}：不确定关系说的是状态本身的宽度，不是扰动；取一大批完全相同制备的电子，一半测位置、一半测动量，两台完美仪器得到的分散度乘积仍有下限。要绕开它，需要的不是更轻的“看”，而是推翻量子态这个概念。
\item 回到袜子抽屉：摸出黑袜子后“剩下是白的”概率从 $\tfrac12$ 变 $1$，和光子过片后的状态更新公式一样。请说出至少两条实质差别，并说明哪一条能被实验直接检验。\emph{提示}：袜子颜色在摸之前就有确定值且顺序无关；量子测量的顺序会改变结果分布——你的三片偏振片就是判决性实验；而“是否存在预先写好一切答案的总表”是第~48~章贝尔不等式的主题。
\item 量子保密通信里，窃听者为什么不能用“极其轻柔地测一下”来避免暴露？\emph{提示}：她必须选提问方向，选错的概率是一半；选错就改写了光子的态，转发也无法复原；双方抽样核对时，约四分之一的失配会把她供出来。轻柔不能改变“提问必须选一个方向”这件事。
\end{puzzles}
